\section{Principios fundamentales de la ingeniería de software}

La ingeniería de software se concibe como una disciplina \textbf{sistemática, disciplinada y cuantificable} porque su objetivo no es solo programar, sino construir productos de software con calidad predecible, costo controlado y capacidad de evolución \cite{sommerville2016}. Esta visión implica que el desarrollo debe apoyarse en métodos definidos, procesos repetibles y métricas que permitan evaluar tanto el producto como el proceso de construcción \cite{pressman2020}.

Un principio central es que el software debe satisfacer tanto los \textbf{requisitos funcionales} como los \textbf{no funcionales} \cite{sommerville2016}. Los requisitos funcionales describen qué debe hacer el sistema, por ejemplo, procesar datos, autenticar usuarios o generar reportes; los no funcionales especifican cómo debe comportarse, incluyendo rendimiento, seguridad, disponibilidad, portabilidad y facilidad de uso \cite{pressman2020}. La ingeniería de software busca equilibrar ambos tipos de requisitos, porque un sistema puede cumplir su función principal y aun así fracasar si no responde adecuadamente a restricciones de calidad o de contexto \cite{sommerville2016}.

La \textbf{calidad del software} es otro principio esencial. No se limita a la ausencia de errores, sino que abarca atributos como \textbf{mantenibilidad, fiabilidad, usabilidad, seguridad y eficiencia} \cite{iso12207}. La mantenibilidad permite corregir y evolucionar el sistema con menor esfuerzo; la fiabilidad garantiza que el software funcione de manera consistente; la usabilidad facilita que los usuarios interactúen con él de forma eficaz; la seguridad protege la información y los recursos; y la eficiencia busca un uso adecuado del tiempo de procesamiento, memoria y otros recursos \cite{iso12207, pressman2020}. En la práctica, estos atributos se evalúan mediante revisiones, pruebas, métricas y estándares de calidad \cite{sommerville2016}.

El \textbf{ciclo de vida del software} organiza el trabajo en etapas sucesivas y relacionadas entre sí. De forma general, este ciclo incluye \textbf{requerimientos, diseño, implementación, pruebas, despliegue y mantenimiento} \cite{pressman2020}. En la fase de requerimientos se identifican las necesidades del usuario y las restricciones del sistema; en el diseño se define la estructura de la solución; en la implementación se traduce el diseño al código; en las pruebas se verifica el cumplimiento de los requisitos; en el despliegue se entrega el sistema al entorno real; y en el mantenimiento se corrigen errores y se incorporan mejoras \cite{sommerville2016}. Esta secuencia no debe entenderse solo como una lista de pasos, sino como una forma de controlar la complejidad y asegurar trazabilidad entre lo que se necesita y lo que finalmente se construye \cite{pressman2020}.

Entre los principios de diseño más importantes se encuentra la \textbf{modularidad}, que consiste en dividir el sistema en partes relativamente independientes \cite{parnas1972}. La modularidad facilita la comprensión, prueba y mantenimiento del software, porque permite trabajar sobre componentes más pequeños y delimitados \cite{parnas1972}. Relacionado con ello está la \textbf{cohesión}, que expresa el grado en que los elementos de un módulo están enfocados a una sola responsabilidad; una cohesión alta suele favorecer la claridad del diseño \cite{pressman2020}. En contraste, el \textbf{acoplamiento bajo} busca minimizar las dependencias entre módulos, de manera que un cambio en una parte del sistema afecte lo menos posible a las demás \cite{sommerville2016}.

Otro principio fundamental es la \textbf{reutilización}, es decir, la posibilidad de emplear componentes, clases, bibliotecas o servicios en diferentes contextos \cite{booch1994}. La reutilización reduce costos, acelera el desarrollo y puede mejorar la confiabilidad cuando los componentes reutilizados ya han sido probados en otros proyectos \cite{gamma1994}. Finalmente, la \textbf{extensibilidad} describe la capacidad del software para incorporar nuevas funciones sin necesidad de reescribir gran parte del sistema \cite{sommerville2016}. Este principio es especialmente valioso en entornos cambiantes, donde los requisitos evolucionan con rapidez y el software debe adaptarse sin perder estabilidad \cite{pressman2020}.

En conjunto, estos principios muestran que la ingeniería de software no consiste únicamente en producir código, sino en construir soluciones correctas, útiles, mantenibles y adaptables. Por ello, la calidad del resultado depende tanto de la gestión del proceso como de las decisiones de diseño adoptadas desde las primeras etapas del desarrollo \cite{sommerville2016, pressman2020}.
